\subsection{Finding a star}
We present here a short proof that every nontrivial class $\calD$ in our setting contains either all cycles, all in-stars, or all out-stars.
\label{sec:findstar}
\begin{proof}[Proof (of Lemma~\ref{lem:findstar})]
  Suppose that there is an integer $\alpha$ such that $\calD$ does not contain a cycle, in-star, or out-star on $\alpha$ vertices. This means in particular that a graph $\calD$ cannot have a strongly connected component with $\alpha$ vertices. Indeed, by selecting $\alpha$ vertices of the component and identifying every other vertex of the graph with one of these vertices, we would obtain a strongly connected graph on $\alpha$ vertices, which is transitively equivalent to a cycle on $\alpha$ vertices.
It also follows that a graph $D\in\calD$ has less than $\alpha$ strongly connected components of size larger than 1: otherwise, identifying one vertex from each such component to a single vertex would create a strongly connected component of size larger than $\alpha$.
  
This bound on the size of the strongly connected components implies that identifying the vertex set of each strongly connected component of some $D\in\calD$ to a single vertex results in an acyclic graph whose size is at most a factor of $\alpha$ smaller. Therefore, the fact that $\calD$ contains arbitrarily large graphs implies that $\calD$ contains arbitrary large acyclic graphs.
Furthermore, $\calD$ contains arbitrarily large graphs without isolated vertices: if we have a graph with $c$ edges and identify every isolated vertex with some vertex $v$, then we get a graph without isolated vertices and exactly $c$ edges. If we identify the strongly connected components of such a graph to single vertices, then less than $\alpha$ isolated vertices can be created (as there are less than $\alpha$ components of size larger than 1). It follows that $\calD$ contains arbitrarily large acyclic graphs with less than $\alpha$ isolated vertices. 

Observe that $D\in\calD$ cannot have a path longer than $\alpha+1$: otherwise, identifying the first and last vertices of the path would create a member of $D\in\calD$ that has a strongly connected component larger than $\alpha$. This means that every acyclic $D\in \calD$ can be partitioned into $\alpha+2$ levels, that is, $D$ has a topological ordering with blocks $B_1$, $\dots$, $B_\ell$ with $\ell\le \alpha+2$ such that each block $B_i$ is an independent set. If an acyclic graph $D_c$ has at least $(\alpha+2)c$ vertices, then one such block $B_i$ has size at least $c$. Let $D^*_c$ be obtained by identifying the blocks $B_1$, $\dots$, $B_{i-1}$ to a single vertex $x$, and the blocks $B_{i+1}$, $\dots$, $B_\ell$ to a single vertex $y$. If there are $\alpha$ vertices in $B_i$ adjacent to both $x$ and $y$, then identifying $x$ and $y$ creates a strongly connected component, a contradiction. If there is a set $I\subseteq B_i$ of  $\alpha$ vertices that is adjacent only to $x$, then identifying $(B_i\setminus I)\cup \{x,y\}$ to a single vertex results in an out-star on $\alpha+1$ vertices. Similarly, if there is a set $I\subseteq B_i$ of  $\alpha$ vertices that is adjacent only to $y$, then identifying $(B_i\setminus I)\cup \{x,y\}$ to a single vertex results in an in-star on $\alpha+1$ vertices. Finally, $I$ cannot have $\alpha$ vertices that are adjacent to neither $x$ or $y$, because they would be isolated vertices. Therefore, if $c\ge 4\alpha$, then we arrive to a contradiction in one of the three cases. 
\end{proof}